% Research integration fragment; no documentclass. Not yet in the main file.
% This supplies an algebraic comparison and a derived local interface.
% It is not an independent verification of Li's Chen constant or a priority claim.

\subsection{The small-initial-cutoff limit of the published four-term sieve}
\label{sec:g1-limit-comparison}

Write $\mu=25/128$, $\Lambda=153/256$, and let
$\Phi(t_1,t_2,t_3)$ be the eight-branch expression in
\href{https://arxiv.org/html/2309.08522v1}{Lichtman,
arXiv:2309.08522v1}, (6.4), with residue exponent $1$ and
spectral exponent $7/32$. In this subsection
\[
 \mathcal T_\mu=\{0<t_3<t_2<t_1<\mu\},\qquad
 S=t_1+t_2+t_3,\qquad d\omega=\frac{dt_1dt_2dt_3}{t_1t_2t_3^2}.
\]
Extend the lower linear-sieve function by $f(s)=0$ for $s\le2$,
and put $\bar f=f/(2e^\gamma)$. Its usual normalization is
$F(s)\to1$, $f(s)\to1$, with $F\ge1\ge f\ge0$.

The published formula (29) in
\href{https://imar.ro/journals/Mathematical_Reports/Pdfs/2026/1_2/3.pdf}
{Runbo Li, \emph{On Chen's theorem, Goldbach's conjecture and almost
prime twins II}} uses the fixed initial exponent $\sigma=1/500$
and gives the numerical upper bound $G_1<6.06932$.
The limiting functional below is an algebraic consequence of that
four-term architecture. It is not a numerical theorem explicitly
printed in that source.

\paragraph{A limit independent of the first three levels.}
For $0<\sigma<\mu$, consider
\begin{align*}
 G_\sigma={}&\frac1\sigma F(\nu_0/\sigma)
 -\frac1\sigma\int_\sigma^\mu
          f((\nu_1(t)-t)/\sigma)\frac{dt}{t}\\
 &+\frac1\sigma\int_{\sigma<t_2<t_1<\mu}
          F((\nu_2(t_1,t_2)-t_1-t_2)/\sigma)
                         \frac{dt_1dt_2}{t_1t_2}\\
 &-\int_{\sigma<t_3<t_2<t_1<\mu}
          f((\Phi(\boldsymbol t)-S)/t_3)\,d\omega.
\end{align*}
Suppose the first three residual exponents have a positive uniform
lower bound:
\[
 \nu_0\ge c,\qquad \nu_1(t)-t\ge c,\qquad
 \nu_2(t_1,t_2)-t_1-t_2\ge c>0.
\]
They may depend on $\sigma$. Then
\begin{equation}\label{eq:g1-limit-identity}
 \lim_{\sigma\downarrow0}G_\sigma
 =\frac1\mu+\int_{\mathcal T_\mu}
       \{1-f((\Phi(\boldsymbol t)-S)/t_3)\}\,d\omega
 =:\mathcal G_\Phi.
\end{equation}
This covers Li's first three levels, since
$\Phi(t_1,t_2,t_3)\ge(1+t_1)/2$ on this domain. It also covers
the more conservative choice of total level $1/2-\delta$ for
\emph{all} of the first three terms: then one may take any
$c<7/64-\delta$, since $1/2-2\mu=7/64$.

Here is a direct convergence proof, avoiding an unquantified
cancellation of divergent integrals. Put $h=F-f$. The
delay equations give $(sh(s))'=-h(s-1)$ for $s>3$.
Thus $h$ is decreasing and
\[
 h'(s)=-\frac{h(s-1)+h(s)}s\le-\frac{2h(s)}s,
 \qquad h(s)\le\frac{9h(3)}{s^2}\quad(s\ge3).
\]
In particular $|F-1|+|1-f|\ll s^{-2}$ suffices here.
The first three terms can therefore be replaced by
\[
 T_\sigma=\frac1\sigma
    \left\{1-\log(\mu/\sigma)+\frac12\log^2(\mu/\sigma)\right\}
\]
with error $O_c(\sigma(1+\log^2(\mu/\sigma)))$.
Direct integration gives
\[
 \int_{\sigma<t_3<t_2<t_1<\mu}d\omega=T_\sigma-\frac1\mu.
\]
Finally
\[
 \Phi-t_1-t_2\ge\frac12-\frac32t_1
       \ge\frac{53}{256}>0.
\]
For sufficiently small $t_3$, the remaining sieve argument is
at least $c_0/t_3$. Consequently the integrand's numerator is
$O(t_3^2)$, and the integrated small-$t_3$ tail is bounded by a
constant times $\int_0^\sigma\log^2(\mu/t_3)\,dt_3$.
This proves both absolute convergence and
\eqref{eq:g1-limit-identity}. The proof only uses the standard
linear-sieve functions and their displayed delay equations;
it does not assume that Li's fixed-$\sigma$ numerical value equals
its limit.

\paragraph{Exact comparison with the compact candidate.}
Take $3\mu<\Theta<3/5$, and set
$L_\Theta(\boldsymbol t)=\min(\Theta,\Phi(\boldsymbol t))$.
The compact first-stop candidate has functional
\begin{equation}\label{eq:compact-g1-functional}
 K_\Theta=
 \frac{2e^{-\gamma}}\mu F(\Theta/\mu)
 +2e^{-\gamma}\int_{\mathcal T_\mu}
       \left\{f((\Theta-S)/t_3)
                    -f((L_\Theta-S)/t_3)\right\}\,d\omega.
\end{equation}
The integrand is zero outside the ordinary third-cube domain
$S+2t_3<\Theta$. Also the ordinary term can be written as
\begin{equation}\label{eq:first-stop-base-identity}
 \frac{2e^{-\gamma}}\mu F(\Theta/\mu)
 =\frac4\Theta+
   4\int_\mu^{\Theta/3}
       \frac{\bar f((\Theta-t)/t)}{t^2}\,dt.
\end{equation}
Indeed, for $3<s<4$,
$sF(s)/(2e^\gamma)=1+\int_2^{s-1}\bar f(u)\,du$.
Thus a first-prefix stop at $\mu$ is simply a change of the
sifting limit to $N^\mu$ at the density-functional level.

Apply the same four-term identity at a constant level $\Theta$.
It follows either by separating an ordinary Rosser upper weight
according to its first zero, one, two, or three primes above
$N^\sigma$, or by the corresponding delay identities. Since
$\Theta>3\mu$, all first-prefix cubes are admitted. Taking
$\sigma\downarrow0$ as above gives
\[
 \frac1\mu F(\Theta/\mu)
 =\frac1\mu+\int_{\mathcal T_\mu}
          \{1-f((\Theta-S)/t_3)\}\,d\omega.
\]
Combining this with \eqref{eq:compact-g1-functional} proves
\begin{equation}\label{eq:g1-cap-exact-comparison}
 K_\Theta=2e^{-\gamma}\mathcal G_\Phi+
 4\int_{\mathcal T_\mu}
       \left\{\bar f((\Phi-S)/t_3)
             -\bar f((\min(\Theta,\Phi)-S)/t_3)\right\}\,d\omega.
\end{equation}
The correction is nonnegative. Every branch of $\Phi$ is at most
$\Lambda=153/256$: the $v$ branches have maximum $\Lambda$, and
each $w$ branch is at most $\Lambda-3t_3/8$ or
$\Lambda-3t_2/8$. Hence equality holds in
\eqref{eq:g1-cap-exact-comparison} whenever $\Theta\ge\Lambda$.
At $\Theta=.597<\Lambda$ the inequality is strict. For example,
near $(t_1,t_2,t_3)=(.11,\mu-.11,.08)$ the ungated
$v(t_1+t_2)$ branch is larger than $.597$, and both lower-sieve
arguments are larger than $2$ on an open subregion.

For the proposed certification values $.596$ and $.597$, all
nonzero third-correction points have
\[
 t_3\ge k_\Theta:=\frac{153/32-8\Theta}{3}>0.
\]
To see this, either $v(t_1)\ge\Theta$, or
$(1+t_1)/2<\Theta$. In the latter case $t_3\le k_\Theta$
makes every inner entry of $w(t_1,t_2,t_3)$ at least $\Theta$;
its gate is also satisfied. Thus $\Phi\ge\Theta$.
The two concrete lower bounds are $53/12000$ and $7/4000$,
respectively. This explains why the capped functional is convenient
for a compact exact enclosure even though the uncapped limit is
numerically a little smaller.

\paragraph{What an exact improvement would mean.}
Li prints a fixed-$\sigma$ bound, not an optimality assertion.
A new certificate $K_\Theta<K_*$ would imply
$2e^{-\gamma}\mathcal G_\Phi<K_*$ by
\eqref{eq:g1-cap-exact-comparison}. Together with the local
distribution interface below, it would give an improved explicit
consequence of the published sieve method. This does not make
that newly certified consequence pre-existing prior art, and it
does not establish a new analytic method or world priority.
An uncomputed possible optimization must not be substituted for
an established numerical comparison bound.

The shifted prime-product proof and the order of limits are written
in the separate fragment \texttt{g1-optimized-transfer.tex}.
